\chapter{Homotopy Equivalence and Contractibility}
\label{ch:viii02-homotopy-equivalence}

Homotopy of maps becomes a relation between spaces when maps are available in
both directions.  The resulting notion, homotopy equivalence, is much weaker
than homeomorphism but strong enough to preserve the invariants that algebraic
topology is designed to compute.  Contractible spaces, deformation retracts,
and homotopy type are the central themes of this chapter.

\section{Homotopy equivalence}

\begin{definition}[Homotopy equivalence]
\label{def:viii02-equivalence}
Spaces \(X\) and \(Y\) are \emph{homotopy equivalent}, written
\[
X\simeq Y,
\]
if there exist continuous maps
\[
f:X\to Y,
\qquad
g:Y\to X
\]
such that
\[
g\circ f\simeq \operatorname{id}_X,
\qquad
f\circ g\simeq \operatorname{id}_Y.
\]
The maps \(f\) and \(g\) are called homotopy inverses.
\end{definition}

\begin{warning}[Not an actual inverse]
\label{warn:viii02-not-inverse}
A homotopy inverse need not satisfy \(g\circ f=\operatorname{id}_X\) or
\(f\circ g=\operatorname{id}_Y\) pointwise.  Equality is replaced by homotopy.
\end{warning}

\section{Homotopy type}

\begin{definition}[Homotopy type]
\label{def:viii02-type}
Two spaces have the same \emph{homotopy type} if they are homotopy equivalent.
\end{definition}

\begin{proposition}[Homotopy equivalence is an equivalence relation]
\label{prop:viii02-equivalence-relation}
Homotopy equivalence is reflexive, symmetric, and transitive.
\end{proposition}

\begin{proof}
Reflexivity uses the identity.  Symmetry is built into the definition.
For transitivity, if \(X\simeq Y\) via \(f,g\) and \(Y\simeq Z\) via \(u,v\),
then \(u\circ f:X\to Z\) and \(g\circ v:Z\to X\) are homotopy inverses.  The
verification uses compatibility of homotopy with composition from VIII/01.
\end{proof}

\section{Contractible spaces}

\begin{definition}[Contractible space]
\label{def:viii02-contractible}
A space \(X\) is \emph{contractible} if
\[
\operatorname{id}_X
\]
is homotopic to a constant map.
\end{definition}

\begin{theorem}[Contractible means homotopy equivalent to a point]
\label{thm:viii02-point}
A nonempty space \(X\) is contractible if and only if
\[
X\simeq \{*\}.
\]
\end{theorem}

\begin{proof}
If \(X\) is contractible to \(x_0\), let \(p:X\to\{*\}\) be the unique map and
\(i:\{*\}\to X\) send the point to \(x_0\).  Then \(p\circ i\) is the identity
on the point, while \(i\circ p\) is the constant map at \(x_0\), homotopic to
\(\operatorname{id}_X\).  The converse is the same argument read backwards.
\end{proof}

\section{Convex and star-shaped examples}

\begin{proposition}[Convex sets are contractible]
\label{prop:viii02-convex}
Every nonempty convex subset \(C\) of a real topological vector space is
contractible.
\end{proposition}

\begin{proof}
Fix \(c_0\in C\) and use
\[
H(x,t)=(1-t)x+tc_0.
\]
\end{proof}

\begin{definition}[Star-shaped set]
\label{def:viii02-star}
A subset \(A\subset\mathbb R^n\) is star-shaped with center \(a\in A\) if the
line segment from \(a\) to every \(x\in A\) lies in \(A\).
\end{definition}

\begin{proposition}[Star-shaped sets are contractible]
\label{prop:viii02-star-contract}
Every nonempty star-shaped subset of \(\mathbb R^n\) is contractible.
\end{proposition}

\begin{example}
Intervals, disks, balls, cubes, and \(\mathbb R^n\) are contractible.
\end{example}

\section{Retractions}

\begin{definition}[Retraction]
\label{def:viii02-retraction}
If \(A\subset X\), a \emph{retraction} is a map
\[
r:X\to A
\]
such that
\[
r|_A=\operatorname{id}_A.
\]
Then \(A\) is called a retract of \(X\).
\end{definition}

Let \(i:A\hookrightarrow X\) denote inclusion.  The retraction condition is
\[
r\circ i=\operatorname{id}_A.
\]

\begin{warning}[Retraction alone is not a deformation retract]
\label{warn:viii02-retract}
A retraction provides an exact one-sided inverse to inclusion, but does not
require \(i\circ r\) to be homotopic to \(\operatorname{id}_X\).
\end{warning}

\section{Deformation retracts}

\begin{definition}[Deformation retract]
\label{def:viii02-deformation-retract}
A subspace \(A\subset X\) is a \emph{deformation retract} of \(X\) if there is a
retraction \(r:X\to A\) such that
\[
i\circ r\simeq \operatorname{id}_X.
\]
Equivalently, there is a homotopy from \(\operatorname{id}_X\) to a map with
image in \(A\), ending at the retraction followed by inclusion.
\end{definition}

\begin{theorem}[Deformation retracts preserve homotopy type]
\label{thm:viii02-dr-equivalence}
If \(A\) is a deformation retract of \(X\), then
\[
A\simeq X.
\]
\end{theorem}

\begin{proof}
The inclusion \(i:A\to X\) and retraction \(r:X\to A\) satisfy
\(r\circ i=\operatorname{id}_A\) and
\(i\circ r\simeq\operatorname{id}_X\).
\end{proof}

\section{Strong deformation retracts}

\begin{definition}[Strong deformation retract]
\label{def:viii02-strong}
A deformation retract \(A\subset X\) is \emph{strong} if the deformation
\[
H:X\times I\to X
\]
from \(\operatorname{id}_X\) to \(i\circ r\) fixes \(A\) pointwise:
\[
H(a,t)=a
\qquad(a\in A,\ t\in I).
\]
\end{definition}

Thus a strong deformation retract is a relative homotopy from VIII/01.

\section{Radial examples}

\begin{example}[Punctured Euclidean space]
\label{ex:viii02-punctured-radial}
The sphere \(S^{n-1}\) is a strong deformation retract of
\(\mathbb R^n\setminus\{0\}\).  Put
\[
r(x)=\frac{x}{\|x\|}
\]
and
\[
H(x,t)=\left((1-t)+\frac{t}{\|x\|}\right)x.
\]
The scalar coefficient is positive, so the homotopy never reaches \(0\).
For \(x\in S^{n-1}\), \(H(x,t)=x\).
\end{example}

\begin{corollary}
\[
\mathbb R^n\setminus\{0\}\simeq S^{n-1}.
\]
\end{corollary}

\begin{example}[Annulus]
\label{ex:viii02-annulus}
For
\[
A=\{x\in\mathbb R^2:1\le\|x\|\le2\},
\]
the unit circle is a strong deformation retract by radial normalization.
\end{example}

\section{Cylinder and circle}

\begin{example}[Cylinder]
\label{ex:viii02-cylinder}
The cylinder \(S^1\times I\) strongly deformation retracts onto
\(S^1\times\{0\}\) by
\[
H((z,s),t)=(z,(1-t)s).
\]
Hence
\[
S^1\times I\simeq S^1.
\]
\end{example}

\section{Graphs with dangling trees}

A contractible branch attached to a graph often contributes no new homotopy
type if it can be collapsed continuously while the remaining graph is fixed.

\begin{example}[Circle with a whisker]
\label{ex:viii02-whisker}
Attach an interval to \(S^1\) at one endpoint.  Collapsing the interval linearly
to its attachment point gives a strong deformation retract onto \(S^1\).
\end{example}

This elementary picture foreshadows cell collapses and CW complexes.

\section{Contractibility versus path connectedness}

\begin{proposition}[Contractible implies path connected]
\label{prop:viii02-contract-path}
Every nonempty contractible space is path connected.
\end{proposition}

\begin{proof}
Let \(H\) contract \(X\) to \(x_0\).  For each \(x\), the path
\(t\mapsto H(x,t)\) joins \(x\) to \(x_0\).  Concatenating such paths joins any
two points.
\end{proof}

\begin{warning}
The converse fails.  The circle \(S^1\) is path connected but is not
contractible.  Proving noncontractibility requires an invariant; degree in
VIII/03 gives one route, and later fundamental-group or homology calculations
give others.
\end{warning}

\section{Homotopy invariants}

\begin{definition}[Homotopy invariant]
\label{def:viii02-invariant}
A construction \(F\) on spaces is a \emph{homotopy invariant} if
\[
X\simeq Y
\quad\Longrightarrow\quad
F(X)\cong F(Y)
\]
in the appropriate category.
\end{definition}

The central strategy of algebraic topology is:
\[
\text{space}
\longmapsto
\text{algebraic invariant},
\]
then use an invariant mismatch to prove that two spaces are not homotopy
equivalent.

\section{Homotopy equivalent need not mean homeomorphic}

\begin{example}[Interval and point]
\label{ex:viii02-interval-point}
The interval \(I\) is contractible, so \(I\simeq\{*\}\), but \(I\) is not
homeomorphic to a point.
\end{example}

\begin{example}[Disk and point]
The closed disk \(D^n\) has the homotopy type of a point while having many more
points and nontrivial local geometry.
\end{example}

\begin{warning}[Do not transfer non-homotopy properties]
\label{warn:viii02-properties}
Homotopy equivalence does not preserve dimension, cardinality, smooth
structure, metric data, or local geometric shape in general.
\end{warning}

\section{A useful criterion}

\begin{proposition}[One-sided strict inverse plus homotopy]
\label{prop:viii02-criterion}
If maps \(f:X\to Y\) and \(g:Y\to X\) satisfy
\[
g\circ f=\operatorname{id}_X
\qquad\text{and}\qquad
f\circ g\simeq\operatorname{id}_Y,
\]
then \(f\) and \(g\) are homotopy inverses.
\end{proposition}

Deformation retracts are the main source of this pattern.

\section{Exercises}

\begin{exercise}\label{exr:viii02-01}
State the two equations defining homotopy inverse maps \(f:X\to Y\) and \(g:Y\to X\).
\end{exercise}
\begin{hint}Compose in both orders.\end{hint}
\begin{solution}
\(g\circ f\simeq\operatorname{id}_X\) and
\(f\circ g\simeq\operatorname{id}_Y\).
\end{solution}

\begin{exercise}\label{exr:viii02-02}
Why is every space homotopy equivalent to itself?
\end{exercise}
\begin{hint}Use the identity map twice.\end{hint}
\begin{solution}
The identity is its own strict inverse, hence also its own homotopy inverse.
\end{solution}

\begin{exercise}\label{exr:viii02-03}
Prove symmetry of homotopy equivalence.
\end{exercise}
\begin{hint}The same pair of maps works in the opposite order.\end{hint}
\begin{solution}
If \(f\) and \(g\) are homotopy inverses, then \(g\) and \(f\) are homotopy
inverses after swapping the roles of \(X\) and \(Y\).
\end{solution}

\begin{exercise}\label{exr:viii02-04}
Define a contractible space.
\end{exercise}
\begin{hint}Contract the identity map.\end{hint}
\begin{solution}
\(X\) is contractible when \(\operatorname{id}_X\) is homotopic to a constant map.
\end{solution}

\begin{exercise}\label{exr:viii02-05}
Show \(\mathbb R^n\) is contractible.
\end{exercise}
\begin{hint}Contract to the origin.\end{hint}
\begin{solution}
Use \(H(x,t)=(1-t)x\).
\end{solution}

\begin{exercise}\label{exr:viii02-06}
Show a convex set is contractible.
\end{exercise}
\begin{hint}Choose any point \(c_0\) in it.\end{hint}
\begin{solution}
The straight-line homotopy \((1-t)x+tc_0\) remains in the convex set.
\end{solution}

\begin{exercise}\label{exr:viii02-07}
Show a star-shaped set is contractible.
\end{exercise}
\begin{hint}Use its star center.\end{hint}
\begin{solution}
The segment from each \(x\) to the center lies in the set, giving the usual
linear contraction.
\end{solution}

\begin{exercise}\label{exr:viii02-08}
Define a retraction \(r:X\to A\).
\end{exercise}
\begin{hint}What must happen on \(A\)?\end{hint}
\begin{solution}
It must satisfy \(r(a)=a\) for all \(a\in A\).
\end{solution}

\begin{exercise}\label{exr:viii02-09}
Translate the retraction condition using the inclusion \(i:A\hookrightarrow X\).
\end{exercise}
\begin{hint}Compose \(r\) after \(i\).\end{hint}
\begin{solution}
\(r\circ i=\operatorname{id}_A\).
\end{solution}

\begin{exercise}\label{exr:viii02-10}
What additional condition turns a retract into a deformation retract?
\end{exercise}
\begin{hint}Compare \(i\circ r\) with the identity of \(X\).\end{hint}
\begin{solution}
One requires \(i\circ r\simeq\operatorname{id}_X\).
\end{solution}

\begin{exercise}\label{exr:viii02-11}
What makes a deformation retract strong?
\end{exercise}
\begin{hint}The subspace should remain fixed during the homotopy.\end{hint}
\begin{solution}
The deformation from \(\operatorname{id}_X\) to \(i\circ r\) must be relative
to \(A\).
\end{solution}

\begin{exercise}\label{exr:viii02-12}
Prove a deformation retract has the same homotopy type as the ambient space.
\end{exercise}
\begin{hint}Use inclusion and retraction as the two maps.\end{hint}
\begin{solution}
They satisfy \(r i=\operatorname{id}_A\) and
\(i r\simeq\operatorname{id}_X\), so they are homotopy inverses.
\end{solution}

\begin{exercise}\label{exr:viii02-13}
Write the radial retraction \(\mathbb R^n\setminus\{0\}\to S^{n-1}\).
\end{exercise}
\begin{hint}Normalize the vector.\end{hint}
\begin{solution}
\(r(x)=x/\|x\|\).
\end{solution}

\begin{exercise}\label{exr:viii02-14}
Why does the radial deformation avoid the origin?
\end{exercise}
\begin{hint}Inspect the scalar multiplying \(x\).\end{hint}
\begin{solution}
The coefficient \((1-t)+t/\|x\|\) is positive, so a nonzero \(x\) remains nonzero.
\end{solution}

\begin{exercise}\label{exr:viii02-15}
Show the radial deformation fixes \(S^{n-1}\).
\end{exercise}
\begin{hint}Set \(\|x\|=1\).\end{hint}
\begin{solution}
The scalar coefficient becomes \((1-t)+t=1\), so \(H(x,t)=x\).
\end{solution}

\begin{exercise}\label{exr:viii02-16}
Show \(S^1\times I\simeq S^1\).
\end{exercise}
\begin{hint}Collapse the interval coordinate.\end{hint}
\begin{solution}
The homotopy \(H((z,s),t)=(z,(1-t)s)\) strongly deformation retracts onto
\(S^1\times\{0\}\cong S^1\).
\end{solution}

\begin{exercise}\label{exr:viii02-17}
Why is a contractible space path connected?
\end{exercise}
\begin{hint}Follow each point along its contraction path.\end{hint}
\begin{solution}
Every point is joined by the contraction to the same terminal point, so any two
points can be joined through that point.
\end{solution}

\begin{exercise}\label{exr:viii02-18}
Give a path-connected space that is not contractible.
\end{exercise}
\begin{hint}The basic example is one-dimensional and closed.\end{hint}
\begin{solution}
\(S^1\).
\end{solution}

\begin{exercise}\label{exr:viii02-19}
Give homotopy-equivalent spaces with different cardinalities.
\end{exercise}
\begin{hint}Use a contractible continuum and a point.\end{hint}
\begin{solution}
The interval \(I\) and a one-point space are homotopy equivalent but have
different cardinalities.
\end{solution}

\begin{exercise}\label{exr:viii02-20}
Does homotopy equivalence preserve dimension?
\end{exercise}
\begin{hint}Compare \(D^n\) with a point.\end{hint}
\begin{solution}
No.  Every disk \(D^n\) is homotopy equivalent to a point.
\end{solution}

\begin{exercise}\label{exr:viii02-21}
If \(X\simeq Y\) and \(Y\simeq Z\), what maps exhibit \(X\simeq Z\)?
\end{exercise}
\begin{hint}Compose the forward maps and compose the backward maps.\end{hint}
\begin{solution}
If \(f:X\to Y,\ g:Y\to X,\ u:Y\to Z,\ v:Z\to Y\), use
\(u f:X\to Z\) and \(g v:Z\to X\).
\end{solution}

\begin{exercise}\label{exr:viii02-22}
Why is a circle with a contractible whisker homotopy equivalent to the circle?
\end{exercise}
\begin{hint}Collapse the whisker to its attachment point.\end{hint}
\begin{solution}
That collapse is a strong deformation retraction onto the circle.
\end{solution}

\begin{exercise}\label{exr:viii02-23}
What is a homotopy invariant?
\end{exercise}
\begin{hint}It cannot distinguish homotopy-equivalent spaces.\end{hint}
\begin{solution}
It is a construction whose value is isomorphic for homotopy-equivalent spaces.
\end{solution}

\begin{exercise}\label{exr:viii02-24}
How can a homotopy invariant prove \(X\not\simeq Y\)?
\end{exercise}
\begin{hint}Compare the invariant values.\end{hint}
\begin{solution}
If the invariant values are not isomorphic, homotopy invariance implies the
spaces cannot be homotopy equivalent.
\end{solution}

\section{Solved problem dossiers}

\begin{problem}[Interval versus point]\label{prob:viii02-01}
Write explicit homotopy-inverse maps between \(I\) and a point.
\end{problem}
\begin{solution}
Let \(p:I\to\{*\}\) be unique and \(i(*)=0\).  Then \(pi\) is the identity of
the point and \(ip\) is the constant map at \(0\), homotopic to
\(\operatorname{id}_I\) via \(H(s,t)=(1-t)s\).
\end{solution}

\begin{problem}[Annulus collapse]\label{prob:viii02-02}
Construct a strong deformation retraction of
\(\{x:1\le\|x\|\le2\}\) onto \(S^1\).
\end{problem}
\begin{solution}
Use
\[
H(x,t)=\left((1-t)+\frac{t}{\|x\|}\right)x.
\]
Its radius is \((1-t)\|x\|+t\), which remains in \([1,2]\), and unit vectors are
fixed.
\end{solution}

\begin{problem}[Punctured plane]\label{prob:viii02-03}
Determine the homotopy type of \(\mathbb R^2\setminus\{0\}\).
\end{problem}
\begin{solution}
Radial normalization gives a strong deformation retraction onto \(S^1\), so
\(\mathbb R^2\setminus\{0\}\simeq S^1\).
\end{solution}

\begin{problem}[Punctured three-space]\label{prob:viii02-04}
Determine the homotopy type of \(\mathbb R^3\setminus\{0\}\).
\end{problem}
\begin{solution}
The same radial construction retracts it strongly onto \(S^2\).
\end{solution}

\begin{problem}[Cylinder core]\label{prob:viii02-05}
Show \(S^1\times[-1,1]\) has the homotopy type of \(S^1\).
\end{problem}
\begin{solution}
Collapse the second coordinate linearly:
\(H((z,s),t)=(z,(1-t)s)\).
\end{solution}

\begin{problem}[Whisker graph]\label{prob:viii02-06}
Let \(X=S^1\cup_{p}I\) attach \(0\in I\) to \(p\in S^1\).  Give a strong
deformation retraction \(X\to S^1\).
\end{problem}
\begin{solution}
Fix every point of \(S^1\), and send a whisker point \(s\in I\) to
\((1-t)s\) along the whisker.  At \(t=1\) the entire whisker is at \(p\).
\end{solution}

\begin{problem}[Product of contractible spaces]\label{prob:viii02-07}
Prove \(X\times Y\) is contractible if \(X\) and \(Y\) are contractible.
\end{problem}
\begin{solution}
Take contractions \(H\) of \(X\) and \(K\) of \(Y\), and define
\(L((x,y),t)=(H(x,t),K(y,t))\).
\end{solution}

\begin{problem}[Homotopy type of a strip]\label{prob:viii02-08}
Determine the homotopy type of \(\mathbb R\times[0,1]\).
\end{problem}
\begin{solution}
The strip is convex, hence contractible and therefore homotopy equivalent to a point.
\end{solution}

\begin{problem}[Retract versus deformation retract]\label{prob:viii02-09}
Explain logically why the equation \(ri=\operatorname{id}_A\) alone does not
establish \(A\simeq X\).
\end{problem}
\begin{solution}
Homotopy equivalence requires information about both composites.  The retraction
equation controls \(ri\), but without \(ir\simeq\operatorname{id}_X\) there is
no homotopy inverse condition in the other direction.
\end{solution}

\begin{problem}[Transporting an invariant]\label{prob:viii02-10}
Suppose \(F\) is a homotopy invariant and \(A\) deformation retracts from \(X\).
What relation holds between \(F(A)\) and \(F(X)\)?
\end{problem}
\begin{solution}
Since \(A\simeq X\), homotopy invariance gives \(F(A)\cong F(X)\).
\end{solution}

\begin{problem}[Noncontractibility strategy]\label{prob:viii02-11}
What would be enough to prove that \(S^1\) is not contractible?
\end{problem}
\begin{solution}
It is enough to find any homotopy invariant that takes different values on
\(S^1\) and on a point.  VIII/03 introduces degree as an early obstruction,
while later chapters supply fundamental group and homology.
\end{solution}

\begin{problem}[Bridge to VIII/03]\label{prob:viii02-12}
Why does the notion of homotopy invariant naturally lead to degree?
\end{problem}
\begin{solution}
To distinguish maps and spaces up to homotopy, one seeks quantities unchanged
by deformation.  Degree assigns an integer to maps between oriented spheres
(or, more generally, oriented closed manifolds) and is invariant under homotopy,
so it supplies the first quantitative obstruction in this volume.
\end{solution}
